\chapter{Unconstrained optimization}

We consider the general problem
\begin{equation}
\min_{x \in \mathcal{R}^n} f(x).
\label{eq:prob}
\end{equation}

\section{Linesearch methods}

\subsubsection{BFGS Linesearch}

\section{Trust-region methods}

\subsection{Unconstrained optimization}

\section{Hessian approximations}

While the gradient $\nabla f(x)$ usually must be known accurately in
minimization algorithms, we often can rely on approximations of the
Hessian $\nabla^2 f(x)$.
This is especially appealing since finite or central differences
require $O(n^2)$ function evaluations, which is prohibitive in many
applications.
In this section, we will focus on traditional secant updates, by
focusing on their definitions, but not on the underlying convergence
theory.
We refer the reader interested in such issues to Chapter 8 of Dennis
and Schnabel~\cite{DennSchn83}.

Reefering to the original problem~\ref{eq:prob}, suppose that we have
just taken a step from $x_c$ to $x_+$.
Following Dennis and Schnabel, we consider the Hessian as the
derivative matrix of the nonlinear equations system $\nabla f(x) =
0$.
The \Index{secant} equation for the unconstrained minimization problem
is
\begin{equation}
H_+s_c = y_c,
\label{eq:secant}
\end{equation}
where
\begin{align*}
s_c &= x_+-x_c,\\
y_c &= \nabla f(x_+) - \nabla f(x_c).
\end{align*}
Since the Hessian of some function is symmetric, we seek an $H_+$ that
satisfies $H_+ = H_+^T$ as well as \eqref{eq:secant}.
The first secant update successfully used in the litterature is the
symmetric secand update\index{secant}

\subsubsection{BFGS}

A popular approximation is the BFGS (Broyden, Fletcher,
Goldfarb and Shanno):
\[
B_{k+1} = B_k + \frac{y_k y_k^T}{y_k^td_k} + \frac{B_kd_k
  (B_kd_k)^T}{d_k^TB_kd_k},
\]
where $y_k = \nabla_x f(x_{k+1}) - \nabla_x f(x_k)$.

\subsubsection{SR1}


\subsection{Trust-region radius update}

\subsubsection{Initial trust-region radius}

Various options again exist to choose the initial trust-region as well
as to update it.

The initialization strategy can be defined by setting the
\verb|radius_init| element of the \verb|BTR| structure to the
desired user's function, or to NULL, if no initialization has to be
performed.
In this case, the initial trust-region radius is affected to
\verb|BTR_INITIAL_DELTA|, defined in the file \verb|defaults.h|.

The default behaviour of the algorithm sets the initial trust-region
to $\|\nabla f(x_0)\|$, as described by Lin and More~\cite{LinMore99},
while more advanced strategies exist. On the first iteration, the
predicted step is used to compute a more reasonable trust-region
radius, using the formula:
\[
\Delta := \min \lbrace \Delta, \| s_0 \| \rbrace.
\]

